\documentclass[12pt]{article}
\usepackage[utf8]{inputenc}
\usepackage[T1]{fontenc}
\usepackage[margin=1in]{geometry}
\usepackage{setspace}
\usepackage{xcolor}
\usepackage[colorlinks=true, urlcolor=blue!60!black, linkcolor=blue!60!black]{hyperref}
\setlength{\parindent}{2em}
\setlength{\parskip}{1.25em}
\renewcommand{\baselinestretch}{1.0}
\title{\textit{Dolan v. City of Tigard}, 512 U.S. 374 (1994)}
\author{Case Brief}
\date{}
\begin{document}
\maketitle

\section{Case Metadata}

\begin{description}
  \item[Case Name:] \emph{Dolan v.\ City of Tigard}
  \item[Citation:] \href{https://supreme.justia.com/cases/federal/us/512/374/}{512 U.S. 374 (1994)}
  \item[Court:] Supreme Court of the United States (on certiorari to the Supreme Court of Oregon, 317 Ore.\ 110, 854 P.2d 437 (1993))
  \item[Argued:] March 23, 1994 \quad \textbf{Decided:} June 24, 1994
  \item[Vote:] 5--4
  \item[Majority Opinion:] Chief Justice William H. Rehnquist (joined by Justices O'Connor, Scalia, Kennedy, and Thomas)
  \item[Dissenting Opinions:] Justice John Paul Stevens (joined by Justices Blackmun and Ginsburg); Justice David H. Souter
  \item[Provisions at Issue:] U.S. Const. amend.~V (Takings Clause), applicable to the States through amend.~XIV; City of Tigard Community Development Code (CDC) \S\S~18.86.040, 18.120.180
  \item[Chapter Placement:] Chapter 6: Substantive Due Process --- Takings \& Property Rights
\end{description}

\section{Facts and Procedural History}

Florence Dolan owned a 9,700-square-foot plumbing and electric supply store on a 1.67-acre parcel on Main Street in the Central Business District of Tigard, Oregon, a city of about 30,000 residents southwest of Portland (\href{https://supreme.justia.com/cases/federal/us/512/374/}{512 U.S. at 377, 379}). The parcel included a gravel parking lot, and Fanno Creek flowed across its southwestern corner and western boundary, placing part of the lot inside the creek's 100-year floodplain. Under Oregon's statewide land-use program, Tigard had codified a Community Development Code (CDC) requiring Central Business District lots to maintain 15\% open space and landscaping, adopting a Master Drainage Plan to preserve the Fanno Creek floodplain as a greenway, and adopting a pedestrian/bicycle pathway plan to relieve downtown traffic congestion (\href{https://supreme.justia.com/cases/federal/us/512/374/}{512 U.S. at 377--379}).

Dolan applied for a redevelopment permit to nearly double her store to 17,600 square feet, pave a 39-space parking lot, and later add a second commercial building (\href{https://supreme.justia.com/cases/federal/us/512/374/}{512 U.S. at 379}). The City Planning Commission approved her permit on two conditions: she had to dedicate to the City (1) all portions of her parcel within Fanno Creek's 100-year floodplain as a public greenway for storm drainage, and (2) an additional 15-foot strip adjacent to the floodplain for a public pedestrian/bicycle pathway (\href{https://supreme.justia.com/cases/federal/us/512/374/}{512 U.S. at 379--380}). Together, the required dedications totaled roughly 7,000 square feet---about 10\% of Dolan's parcel---which could count toward the CDC's 15\% open-space requirement.

The Commission denied Dolan's variance request, finding that increased impervious surfaces ``can only add to the public need to manage the stream channel and floodplain,'' and that a pathway ``could offset some of the traffic demand'' generated by the larger store (\href{https://supreme.justia.com/cases/federal/us/512/374/}{512 U.S. at 381--382}). After the Tigard City Council approved the order, Dolan appealed to the Oregon Land Use Board of Appeals (LUBA), arguing that the dedications were an uncompensated taking under the Fifth Amendment. LUBA, the Oregon Court of Appeals (113 Ore.\ App.\ 162), and the Oregon Supreme Court (317 Ore.\ 110) all affirmed, concluding that the conditions bore a ``reasonable relationship'' to the expansion's impacts (\href{https://supreme.justia.com/cases/federal/us/512/374/}{512 U.S. at 382--383}). The Supreme Court granted certiorari and reversed.

\section{Issue(s) Presented}

Building on \href{https://supreme.justia.com/cases/federal/us/483/825/}{\emph{Nollan v.\ California Coastal Commission}, 483 U.S. 825 (1987)}, what degree of connection does the Fifth Amendment Takings Clause require between land-dedication conditions imposed on a development permit and the projected impacts of the proposed development, and did the City of Tigard's findings satisfy that standard?

\section{Holding}

\textbf{Reversed and remanded (5--4).} Even when a permit condition satisfies \emph{Nollan}'s ``essential nexus'' requirement, it violates the Takings Clause unless the government makes an individualized determination demonstrating ``rough proportionality''---that the required dedication is related both in nature and extent to the projected impact of the proposed development (\href{https://supreme.justia.com/cases/federal/us/512/374/}{512 U.S. at 386, 391}). Tigard failed to show why flood control required a \emph{public} greenway rather than private open space, and failed to quantify that the pedestrian/bicycle pathway would---rather than merely ``could''---offset the store's added traffic (\href{https://supreme.justia.com/cases/federal/us/512/374/}{512 U.S. at 392--396}).

\section{Reasoning}

Chief Justice Rehnquist observed that had Tigard simply demanded a public strip of land along Fanno Creek outright, it would unquestionably have committed a per se taking by stripping Dolan of her right to exclude others (\href{https://supreme.justia.com/cases/federal/us/512/374/}{512 U.S. at 384}). Unlike legislative, citywide zoning rules upheld under \emph{Village of Euclid v.\ Ambler Realty Co.}, 272 U.S. 365 (1926), Tigard made an adjudicative decision conditioning an individual building permit on deeding land to the City. Under the ``unconstitutional conditions'' doctrine, the government cannot force an owner to surrender her Fifth Amendment right to just compensation in exchange for a discretionary permit where the exaction lacks the required relationship to the development's impacts (\href{https://supreme.justia.com/cases/federal/us/512/374/}{512 U.S. at 385}).

\begin{itemize}
  \item \textbf{Step 1 --- Essential Nexus (\emph{Nollan}) Satisfied:} Unlike \href{https://supreme.justia.com/cases/federal/us/483/825/}{\emph{Nollan}}, where a lateral beach easement had no logical tie to preserving ocean views, Tigard's conditions met Step 1: preventing flooding along Fanno Creek and reducing downtown traffic congestion are legitimate public purposes, and limiting floodplain development and providing alternative bike/pedestrian routes logically relate to those ends (\href{https://supreme.justia.com/cases/federal/us/512/374/}{512 U.S. at 386--388}).
  \item \textbf{Step 2 --- ``Rough Proportionality'' Standard and Burden of Proof:} Surveying state exaction cases, the Court rejected both overly lax ``generalized statement'' tests and Illinois's overly rigid ``specifically and uniquely attributable'' test, endorsing instead an intermediate standard akin to the majority state-court ``reasonable relationship'' test (\href{https://supreme.justia.com/cases/federal/us/512/374/}{512 U.S. at 389--391}). To avoid confusion with Equal Protection ``rational basis'' review, the Court coined the term ``rough proportionality'': no precise mathematical calculation is required, but the \emph{city} bears the burden of making an individualized determination that the dedication is related in both nature and extent to the development's impact (\href{https://supreme.justia.com/cases/federal/us/512/374/}{512 U.S. at 391 \& n.8}).
  \item \textbf{Application to Floodplain Greenway (Failed):} Keeping the floodplain undeveloped served flood control, and Tigard's 15\% open-space rule already achieved that end. Tigard never explained why a \emph{public} greenway---eviscerating Dolan's right to exclude recreational visitors---was necessary for flood control (\href{https://supreme.justia.com/cases/federal/us/512/374/}{512 U.S. at 392--395}). Distinguishing \emph{PruneYard Shopping Center v.\ Robins}, 447 U.S. 74 (1980), where a mall could impose time, place, and manner limits on leafleters, the Court noted that a permanent public easement stripped Dolan of all control over public entry (\href{https://supreme.justia.com/cases/federal/us/512/374/}{512 U.S. at 394}).
  \item \textbf{Application to Pedestrian/Bicycle Pathway (Failed):} Although Tigard estimated the expansion would generate roughly 435 additional vehicle trips per day, its conclusory finding that a bike path ``could offset some of the traffic demand'' fell short of showing that the path \emph{will}, or is likely to, offset traffic demand in proportion to the 15-foot dedication (\href{https://supreme.justia.com/cases/federal/us/512/374/}{512 U.S. at 395--396}).
\end{itemize}

\section{Concurring and Dissenting Opinions}

No separate concurring opinions were filed. Two dissents challenged both the Court's standard and its burden-shifting framework:

\begin{itemize}
  \item \textbf{Justice Stevens (Dissenting, joined by Blackmun and Ginsburg):} Argued that state-court precedents did not support ``rough proportionality'' or shifting the burden of proof onto municipalities, and faulted the majority for isolating one ``strand'' (the right to exclude) rather than evaluating the parcel as a whole and the reciprocal benefits Dolan gained from the expansion permit and flood-channel improvements (\href{https://supreme.justia.com/cases/federal/us/512/374/}{512 U.S. at 397--404}). Dismissing the ``could'' vs.\ ``would'' critique as a ``play on words,'' Stevens warned that applying heightened unconstitutional-conditions scrutiny to commercial land-use permits resurrected \emph{Lochner}-era substantive due process (\href{https://supreme.justia.com/cases/federal/us/512/374/}{512 U.S. at 405--411}).
  \item \textbf{Justice Souter (Dissenting):} Contended that the majority did not actually apply rough proportionality to these facts, because both defects it identified---no link between public recreation and flood control, and a tenuous ``could offset'' traffic finding---were really \emph{Nollan} nexus rulings driven by improperly shifting the evidentiary burden away from the challenger (\href{https://supreme.justia.com/cases/federal/us/512/374/}{512 U.S. at 411--414}). Souter also noted that Tigard's conditions were imposed pursuant to a pre-existing municipal code rather than ad hoc ``reverse spot'' zoning (\href{https://supreme.justia.com/cases/federal/us/512/374/}{512 U.S. at 413 n.*}).
\end{itemize}

\section{Key Quotations}

\begin{enumerate}
  \item \textbf{Chief Justice Rehnquist on ``Rough Proportionality''} (\href{https://supreme.justia.com/cases/federal/us/512/374/}{512 U.S. at 391}):
  \begin{quote}
    ``We think a term such as `rough proportionality' best encapsulates what we hold to be the requirement of the Fifth Amendment. No precise mathematical calculation is required, but the city must make some sort of individualized determination that the required dedication is related both in nature and extent to the impact of the proposed development.''
  \end{quote}

  \item \textbf{Chief Justice Rehnquist on the Unconstitutional Conditions Doctrine} (\href{https://supreme.justia.com/cases/federal/us/512/374/}{512 U.S. at 385}):
  \begin{quote}
    ``Under the well-settled doctrine of `unconstitutional conditions,' the government may not require a person to give up a constitutional right---here the right to receive just compensation when property is taken for a public use---in exchange for a discretionary benefit conferred by the government where the benefit sought has little or no relationship to the property.''
  \end{quote}

  \item \textbf{Chief Justice Rehnquist on the Status of the Takings Clause} (\href{https://supreme.justia.com/cases/federal/us/512/374/}{512 U.S. at 392}):
  \begin{quote}
    ``We see no reason why the Takings Clause of the Fifth Amendment, as much a part of the Bill of Rights as the First Amendment or Fourth Amendment, should be relegated to the status of a poor relation in these comparable circumstances.''
  \end{quote}

  \item \textbf{Justice Stevens (Dissenting) on Burden-Shifting and Substantive Due Process} (\href{https://supreme.justia.com/cases/federal/us/512/374/}{512 U.S. at 405}):
  \begin{quote}
    ``The Court has made a serious error by abandoning the traditional presumption of constitutionality and imposing a novel burden of proof on a city implementing an admittedly valid comprehensive land use plan. Even more consequential than its incorrect disposition of this case, however, is the Court's resurrection of a species of substantive due process analysis that it firmly rejected decades ago.''
  \end{quote}
\end{enumerate}

\section{Doctrinal Significance and Modern Legacy (Through 2026)}

Together with \href{https://supreme.justia.com/cases/federal/us/483/825/}{\emph{Nollan}}, \emph{Dolan} established the ``\emph{Nollan/Dolan} exactions test,'' a distinct pillar of Takings jurisprudence alongside physical takings (\emph{Loretto}), total wipeouts (\emph{Lucas}), and ad hoc balancing (\emph{Penn Central}) (\href{https://www.law.cornell.edu/wex/takings}{Cornell LII}). Subsequent Supreme Court decisions have clarified and dramatically expanded \emph{Dolan}'s reach:

\begin{itemize}
  \item \textbf{Separating Exactions from Substantive Due Process (\emph{Lingle}):} In \href{https://supreme.justia.com/cases/federal/us/544/528/}{\emph{Lingle v.\ Chevron U.S.A. Inc.}, 544 U.S. 528 (2005)}, the Court repudiated \emph{Agins}'s ``substantially advances'' test for general regulatory takings as a disguised due-process inquiry, but expressly preserved \emph{Nollan} and \emph{Dolan} as a specialized application of the unconstitutional-conditions doctrine.
  \item \textbf{Monetary Exactions and Permit Denials (\emph{Koontz}):} In \href{https://supreme.justia.com/cases/federal/us/570/595/}{\emph{Koontz v.\ St.\ Johns River Water Management District}, 570 U.S. 595 (2013)}, the Court extended \emph{Nollan/Dolan} scrutiny in two major ways: it applies both when a government \emph{denies} a permit because an owner refuses an unconstitutional condition and when the government demands \emph{money} (such as funding off-site mitigation) rather than a deed to real property, while leaving ordinary taxes outside \emph{Dolan}'s scope.
  \item \textbf{Physical Access and Inspection Conditions (\emph{Cedar Point}):} In \href{https://supreme.justia.com/cases/federal/us/594/20-107/}{\emph{Cedar Point Nursery v.\ Hassid}, 594 U.S. \_\_\_ (2021)}, the Court held that regulations granting union organizers intermittent access to farmland effect per se physical takings, while noting that reasonable health-and-safety inspection conditions attached to permits generally satisfy \emph{Nollan/Dolan}.
  \item \textbf{Extension to Legislative Exactions (\emph{Sheetz}, 2024--2026):} For thirty years, lower courts split over whether \emph{Dolan}'s ``adjudicative'' framing exempted legislatively enacted impact-fee schedules. In \href{https://supreme.justia.com/cases/federal/us/601/22-1074/}{\emph{Sheetz v.\ County of El Dorado}, 601 U.S. \_\_\_ (2024)}, a unanimous Court (opinion by Justice Barrett) held that the Takings Clause ``does not distinguish between legislative and administrative permit conditions,'' subjecting El Dorado County's legislatively enacted traffic-impact fee (nearly \$24,000 on a prefabricated home) to \emph{Nollan/Dolan} review (\href{https://www.law.cornell.edu/supremecourt/text/22-1074}{Cornell LII}). However, the Court left open whether class-based legislative fee schedules must be tailored as specifically as parcel-by-parcel exactions---a question Justice Kavanaugh's concurrence emphasized remains unresolved. On June 15, 2026, the Court denied certiorari in Sheetz's post-remand petition (\href{https://www.scotusblog.com/cases/sheetz-v-el-dorado-california/}{\emph{Sheetz v.\ El Dorado}, No.\ 25-958}), leaving lower courts to work out the degree of tailoring required for legislative fee formulas.
\end{itemize}

\section{Pedagogical Critiques and Exam Hypotheticals}

\subsection{Pedagogical Critiques}
\begin{itemize}
  \item \textbf{Burden Inversion and Chilling Effect:} By shifting the burden of proof onto local governments and demanding quantitative findings, \emph{Dolan}---especially after \emph{Koontz} and \emph{Sheetz}---may incentivize municipalities either to deny development permits outright (where lenient \emph{Penn Central} review applies) or to rely on broad property taxes rather than targeted mitigation fees.
  \item \textbf{Unresolved Tailoring for Legislative Fees:} While \emph{Sheetz} closed the legislative-exemption loophole, it created acute doctrinal tension with \emph{Dolan}'s call for an ``individualized determination'': how much aggregation is permissible in a citywide impact-fee formula?
\end{itemize}

\subsection{Exam Hypothetical and ``Opposite Side'' Contrast}
\textbf{Hypothetical:} A coastal city enacts an ordinance requiring any commercial developer seeking a building permit to expand floor space by more than 5,000 square feet to (a) deed a 20-foot public waterfront boardwalk easement across its lot, and (b) pay a flat \$250,000 ``Downtown Cultural Arts Fee.'' A warehouse owner applies to add 6,000 square feet of storage space (adding only 15 truck trips per week) and challenges both conditions under the Fifth Amendment.

\textbf{Analysis:} Under \href{https://supreme.justia.com/cases/federal/us/601/22-1074/}{\emph{Sheetz}} and \href{https://supreme.justia.com/cases/federal/us/570/595/}{\emph{Koontz}}, neither condition escapes heightened scrutiny merely because it is legislatively enacted or demands money. The \$250,000 Cultural Arts Fee fails \href{https://supreme.justia.com/cases/federal/us/483/825/}{\emph{Nollan}}'s Step~1 ``essential nexus'' test because warehouse storage has no logical connection to funding municipal arts. Even if the city argues the waterfront easement offsets coastal congestion, it fails \href{https://supreme.justia.com/cases/federal/us/512/374/}{\emph{Dolan}}'s Step~2 ``rough proportionality'' test: the city has made no individualized showing that 15 weekly truck trips justify stripping the owner of her right to exclude the public from a permanent 20-foot boardwalk.

\textbf{``Opposite Side'' Contrast:} Now alter the conditions: suppose the city instead (1) denies the expansion permit outright under a neutral zoning cap limiting warehouse height and impervious coverage, or (2) conditions the permit on the owner retaining a \emph{private} vegetated stormwater buffer (with no public access) and paying a \$12,000 road-maintenance fee calculated from an engineering study of pavement wear caused by 15 heavy trucks per week. In scenario (1), no exaction is demanded, so \emph{Nollan/Dolan} does not apply and the zoning cap is evaluated under deferential \emph{Penn Central} balancing. In scenario (2), the private stormwater buffer preserves the owner's right to exclude (\href{https://supreme.justia.com/cases/federal/us/512/374/}{\emph{Dolan}, 512 U.S. at 393}), and the calibrated \$12,000 road fee satisfies both \emph{Nollan}'s essential nexus and \emph{Dolan}'s rough proportionality.

\end{document}
